\lecture{3}{2026-09-14}{}

\begin{example}[Average number of zeroes in binary strings]
    Recall that the bivariate generating function for binary strings parameterized by the number of zeroes is
    \[
        B(u, z) = \frac{1}{1 - (uz + z)} = \frac{1}{1 - z(1+u)}.
    \]
    We compute
    \[
        [z^n] B(1, z) = [z^n] \frac{1}{1 - 2z} = 2^n
    \]
    and
    \[
        [z^n] B_u(1, z) = [z^n] \left. \frac{z}{(1 - (1+u)z)^2} \right|_{u=1} = [z^n] \frac{z}{(1 - 2z)^2} = n 2^{n-1},
    \]
    using the generalized binomial theorem.
    Therefore, the expected number of zeroes in a binary string of length \(n\) is
    \[
        \mathbb{E}[\chi_n] = \frac{[z^n] B_u(1, z)}{[z^n] B(1, z)} = \frac{n 2^{n-1}}{2^n} = \frac{n}{2}.
    \]
\end{example}

\begin{example}[Average number of ones in a composition]
    Recall that the unparameterized class \(\mathcal{C}\) of compositions has a specification
    \[
        \mathcal{C} = \operatorname{SEQ}(\mathcal{Z} \times \operatorname{SEQ}(\mathcal{Z})) = \operatorname{SEQ}(\mathcal{Z} + \mathcal{Z}^2 \times \operatorname{SEQ}(\mathcal{Z})).
    \]
    We track the number of ones by replacing the atom \(\mathcal{Z}\) with \(\mathcal{Z} \times \mu\):
    \[
        \mathcal{C} = \operatorname{SEQ}((\mathcal{Z} \times \mu) + \mathcal{Z}^2 \times \operatorname{SEQ}(\mathcal{Z})),
    \]
    which gives the bivariate generating function
    \[
        C(u, z) = \frac{1}{1 - (uz + z^2 \frac{1}{1-z})} = \frac{1-z}{(1-z)(1-uz) - z^2}.
    \]
\end{example}

\section*{Basics of Asymptotics}

Let \(f_n\) and \(g_n\) be functions from \(\mathbb{N}\) to \(\mathbb{R}_{\ge 0}\).
We write:
\begin{itemize}
    \item \(f_n \sim g_n\) if \(\lim_{n\to\infty} f_n/g_n = 1\);
    \item \(f_n = O(g_n)\) if there exists a constant \(C > 0\) and \(N \in \mathbb{N}\) such that
        \[ 
            f_n \le C g_n \quad \text{for all } n \ge N;
        \]
    \item \(f_n = o(g_n)\) if \(\lim_{n\to\infty} f_n/g_n = 0\).
\end{itemize}

Here is a hierarchy of asymptotic growth rates:
\begin{center}
    \begin{tabular}{rl}
        \(A\) & constant \\
        \((\log n)^B\) & polylogarithmic \\
        \(n^C\) & polynomial \\
        \(D^n\) & exponential \\
        \(n^{En}\) & ``\(n\) to the \(n\)''
    \end{tabular}
\end{center}

\section*{Taylor's Theorem}

\begin{theorem}[Taylor's theorem]
    Let \(f\) be a complex differentiable function on the set
    \[
        D = \{z \in \mathbb{C} : |z - \omega| < r\}
    \]
    for some \(\omega \in \mathbb{C}\) and \(r > 0\).
    Then, at \(z = \omega\), \(f\) has a Taylor series expansion
    \[
        f(z) = \sum_{n \ge 0}^k \frac{f^{(n)}(\omega)}{n!} (z - \omega)^n + O\left((z - \omega)^{k+1} e(z)\right),
    \]
    for some complex funciton \(e(z)\) satisfying
    \[
        |e(z)| \le \frac{\max_{|z - \omega| = r} |f(z)|}{r^k(r-|z - \omega|)}.
    \] 
\end{theorem}

\begin{example}
    We know \(\frac{1}{1+z} = \sum_{k \geq 0} (-1)^k z^k\) for \(|z| < 1\).
    Therefore, for \(n > 1\), we have
    \[
        \frac{1}{1+n} = \frac{1}{n} \cdot \frac{1}{1 + \frac{1}{n}} = \frac{1}{n} \sum_{k \geq 0} (-1)^k n^{-k} = \sum_{k \geq 0} (-1)^k n^{-(k+1)}.
    \]
\end{example}



For every \(n \in \mathbb{N}\), let \(S(n)\) be a finite subset of \(\mathbb{N}\); usually \(S(n) = \{0, 1, \dots, n\}\).
How do we find the asymptotics of 
\[
    s_n = \sum_{k \in S(n)} a_n(k),
\]
when \(a_n(k)\) is known explicitly?

Note that its not enough to know the asymptotics of \(a_n(k)\) for each fixed \(k\).

\begin{example}
    Let 
    \[
        a_n(k) = 2^{-k} + 2^{-n+k} \qquad \text{ and } \qquad b_n(k) = 2^{-k}.
    \]
    Then, for each fixed \(k\), we have \(a_n(k) \sim b_n(k)\) as \(n \to \infty\).
    However, we have
    \[
        \sum_{k=0}^n a_n(k) = \sum_{k=0}^n 2^{-k} + \sum_{k=0}^n 2^{-n+k} = 2 - 2^{-n} + 2 - 2^{-n} = 4 - 2^{1-n} \sim 4,
    \]
    while
    \[
        \sum_{k=0}^n b_n(k) = \sum_{k=0}^n 2^{-k} = 2 - 2^{-n} \sim 2.
    \]
\end{example}

We say that \(a_n(k) \sim b_n(k)\) \emph{uniformly} for \(k \in S(n)\) if
\[
    \sup_{k \in S(n)} \left| \frac{a_n(k)}{b_n(k)} - 1 \right| \to 0 \quad \text{as } n \to \infty.
\]

Note: Integrals are limits of sums.
For example, how can we approximate \(\sum_{k=0}^n k^2\)?
We can note
\[
    \frac{n^3}{3}
    =
    \int_0^n x^2 dx
    \leq
    \sum_{k=0}^n k^2
    \leq
    \int_0^{n+1} x^2 dx
    =
    \frac{(n+1)^3}{3}
    \sim
    \frac{n^3}{3}.
\]
Therefore, \(\sum_{k=0}^n k^2 \sim \frac{n^3}{3}\).